\documentclass[11pt]{article}

\usepackage[margin=25mm]{geometry}

\usepackage{amsmath,amssymb,fontspec,unicode-math}

\IfFontExistsTF{DejaVu Serif}{\setmainfont{DejaVu Serif}}{\setmainfont{Latin Modern Roman}}

\IfFontExistsTF{Latin Modern Math}{\setmathfont{Latin Modern Math}}{\setmathfont{DejaVu Math TeX Gyre}}

\usepackage{xurl,hyperref}

\setlength{\parindent}{0pt}\setlength{\parskip}{6pt}

\begin{document}

\title{Smooth Regularization for Efficient Video Recognition}

\author{ELVIA preliminary AI analysis}\date{21 October 2026}\maketitle

Preliminary assessment: Minor.\par

Independent scientific review remains pending. Author adoption has not been requested.

The documented Kinetics test subset differs from the cited MoViNet validation protocol, so unmatched score differences do not isolate model gains. An anchored relative scale theorem also does not imply unit-sphere embeddings. The loss and original appendix theorem survive the scoped checks.

\section*{Review scope}

Claim-focused reading of the original paper and relevant author-released source, with direct mathematical and static-source checks. No author program, historical verifier, model training or published benchmark was executed.

\section*{Problem 1: Cross-model gains require a common population}

Source: Kinetics-600 Tables 1–2; author README, lines 146–156; official pre-paper MoViNet evaluation configuration

Assessment: supported; impact: minor

\subsection*{Source claim (paraphrase)}

Author documentation explicitly identifies still-accessible labeled test videos, while the cited baseline configuration uses validation data. Baseline values match the official model table.

\subsection*{Analysis}

For author set \(S_G\), baseline set \(S_B\), the reported difference decomposes as \[A_G(S_G)-A_B(S_B)=[A_G(S_G)-A_B(S_G)]+[A_B(S_G)-A_B(S_B)].\] The second term is unmeasured. A baseline scoring \(72/100\) on one set and \(78/100\) on another can appear \(6\) points behind an identical-prediction modified model tested only on the latter. This establishes non-identification, not actual video-set difficulty. An unseen matched rerun could resolve it.

\subsection*{Finding and scope}

Matched manifests/predictions are required. The internal attention-only versus GRW ablation may remain informative if its protocol is shared; the strongest split evidence concerns MoViNet, not automatically every row or Kinetics-400.

\subsection*{Proposed correction}

Evaluate all checkpoints on one frozen labeled manifest, with clip/crop and compute conventions explicit. The existing author subset clarification is credited but does not supply matched baseline results.

The authors already published the test-subset clarification. A matched baseline evaluation resolving this comparison is not established.

\section*{Problem 2: A relative span theorem does not imply unit norm}

Source: Affine-normalization discussion; Appendix A, anchored scalar theorem and Propositions 1–2

Assessment: supported; impact: minor

\subsection*{Source claim (paraphrase)}

The main text attributes near-unit-sphere embeddings to affine-map/GRW optimization; the appendix proves a relative anchored bound.

\subsection*{Analysis}

Translation by \(c\) leaves all velocities, accelerations and permuted energies unchanged, so \(L_\alpha(Z+c)=L_\alpha(Z)\). Changing affine bias realizes arbitrarily large norms with identical regularizer. This does not refute the theorem’s anchor \(z_1=0\), or analyze classifier/weight-decay interactions. The anchored scalar span proof remains valid.

\subsection*{Finding and scope}

Relative scale does not establish unit radius or shell concentration; empirical norm claims need their own distribution and normalization convention.

\subsection*{Proposed correction}

For any loss sublevel \(L_\alpha(Z)\le\log((T-1)!)\), nonnegative ordering loss and Cauchy–Schwarz give \[\|z_j-z_1\|\le\sqrt{2(T-1)\log((T-1)!)/\alpha}\]. This is relative and conditional, not a unit-sphere theorem.

A proposed correction is described; author adoption is not established.

\section*{Limits}

Preliminary AI assessment; independent human scientific adjudication remains pending.\par

Analytic examples and release discrepancies do not remeasure the paper’s reported experiments.\par

\section*{Source identity}

arXiv:2511.20928v2 (latest author version confirmed from current official metadata)

\url{https://arxiv.org/pdf/2511.20928v2}

Reviewed PDF SHA-256: \nolinkurl{c59285f47c8d652476681919da9d14394bfed8f2f13a0f67c05f4ac1d668837c}

The SHA-256 identifies the exact retained PDF inspected in this rapid source/argument check. Original supplied report bytes are retained locally. No unambiguous original PDF retrieval timestamp was recovered; file modification time was not substituted. Current official metadata identifies the latest author version selected here; the relevant passages were read in that version.

\end{document}
